% Appendix: Distribution Shift Analysis - Extended Details
% Supplementary material for Section \ref{sec:distribution_shift}

\subsection{Distribution Shift: Extended Analysis}
\label{app:distribution_shift_details}

This appendix provides additional detail for the distribution shift analysis in Section~\ref{sec:distribution_shift}, including robustness checks, PCA methodology justification, configuration guidance for practitioners, and comparison to related work.

\subsubsection{Robustness Analysis Across PC Dimensions}

A natural concern is that PC1 explains only 3.1\% of total variance. To address this, we compute average per-component PSI (quantile-based bins) across different dimensionality choices:

\begin{itemize}[leftmargin=*,noitemsep,topsep=2pt]
    \item PC1 alone: PSI $= 0.225$ (significant)
    \item PC1--5: avg PSI $= 0.197$ (moderate--significant)
    \item PC1--10: avg PSI $= 0.133$ (moderate)
    \item All 33 routing features: avg PSI $= 0.089$ (no significant shift)
\end{itemize}

\noindent The attenuation pattern is expected: leading PCs capture between-domain variation (where shift is concentrated), while higher PCs capture within-domain variation that is more stable across sources. PC1's elevated shift confirms it is the primary axis where prior and deployment distributions differ, validating its role as a routing feature. The fact that average shift across all features remains at moderate levels ($\sim 0.09$) indicates the mismatch is genuine, not a PC1 artifact.

\subsubsection{Why the Production PCA (Not a Generic Baseline)?}

The PCA artifact shipped with the library was trained on 80K RouteLLM battle prompts---a dataset independent of the evaluation prompts (Table~\ref{tab:dataset}). One might ask whether using this domain-adapted PCA inflates the measured shift. The answer is no: because the production router projects \emph{every} incoming prompt through this same PCA, the measured shift is exactly the covariate shift that LinUCB experiences in deployment. Using a generic PCA (e.g., trained on C4 web text) would measure shift in a feature space the router never sees, answering a different and less actionable question. If anything, the production PCA may \emph{underestimate} total shift, because it projects away variance directions present in deployment data but absent from its RouteLLM training set. Figure~\ref{fig:lmsys_holdout_structure} separately validates that this PCA generalizes across data sources (Spearman $\rho = -0.370$, $p < 0.0001$, exceeding all 100 random projections).

\subsubsection{Configuration Guidance for Library Users}

The distribution shift analysis translates into concrete configuration recommendations:

\begin{itemize}[leftmargin=*,noitemsep,topsep=2pt]
    \item \textbf{Use warm-start priors for Day~1 intelligence.} Despite the shift, the shipped priors achieve 28\% lower regret than cold-start on the holdout (Table~\ref{tab:performance_gap}). The shift is significant but small in practical terms (d $= 0.33$), so the priors remain a valuable initialization.
    \item \textbf{Enable online learning (the default).} Leave the bandit's update mechanism active so it corrects for prior miscalibration as deployment data accumulates. The library's default configuration (\texttt{priors="warmup"}) already does this.
    \item \textbf{Enable Corralling for safety-critical deployments.} If routing errors have high cost (e.g., customer-facing applications with SLA requirements), enable Corralling to bound worst-case regret during the adaptation period.
    \item \textbf{Monitor distribution metrics.} The library exposes PSI and feature statistics via \texttt{router.stats()}. Tracking PSI over time allows users to detect drift and, if necessary, trigger prior regeneration (\texttt{generate\_warmup\_priors.py}).
\end{itemize}

\subsubsection{Comparison to Related Work}

Prior work on LLM routing~\cite{ong2024routellm,chen2024frugalgpt} typically assumes the deployment distribution matches the training distribution---an assumption our analysis shows is violated even within LLM evaluation domains. RouteLLM~\cite{ong2024routellm} trains matrix factorization models on battle data but provides no mechanism to detect or adapt to distribution shift; if a user deploys RouteLLM on a prompt population that differs from its training data, routing quality degrades silently. FrugalGPT~\cite{chen2024frugalgpt} uses fixed cascading policies that cannot update at all. In contrast, BanditGPT:

\begin{enumerate}[leftmargin=*,noitemsep,topsep=2pt]
    \item \textbf{Diagnoses shift:} The PSI analysis (0.225, Cohen's d $= 0.33$) provides the first quantitative evidence that distribution shift is a practical concern for LLM routing, not merely a theoretical risk
    \item \textbf{Adapts continuously:} LinUCB updates correct for prior miscalibration as deployment data arrives, requiring no manual intervention
    \item \textbf{Provides safety guarantees:} Corralling bounds worst-case regret via expert hedging (Table~\ref{tab:performance_gap}), a property no other LLM router offers
\end{enumerate}

\noindent The task-category decomposition (Figure~\ref{fig:distribution_shift}, bottom) further validates the library's feature engineering: despite extensive overlap (Cohen's d $= 0.23$), the PCA-reduced features carry a routing-relevant signal that distinguishes Mixtral-Sufficient from GPT-4-Turbo-Required prompts---exactly the signal the library's bandit exploits over many rounds.

\subsubsection{Detailed Statistical Results}

\begin{itemize}[leftmargin=*,noitemsep,topsep=2pt]
    \item \textbf{PSI:} 0.225 (95\% bootstrap CI: [0.194, 0.285]; 1000 resamples, quantile-based bins with $B=20$)
    \item \textbf{Cohen's d:} 0.33 (small effect, comparing PC1 distributions)
    \item \textbf{KS test:} D $= 0.119$, $p < 10^{-22}$ (distributional difference confirmed)
    \item \textbf{Mean PC1 shift:} $+0.060$ (deployment prompts shifted rightward on PC1)
    \item \textbf{Task separation (within prior):} Cohen's d $= 0.23$ between Mixtral-Sufficient and GPT-4-Turbo-Required, with $>$99\% range overlap
    \item \textbf{Reference distribution:} Complete 80K RouteLLM prior
    \item \textbf{Deployment distribution:} 1,871 unique LMSYS prompts (1,121 dev $+$ 750 holdout)
\end{itemize}
